\chapter{Универсальное покрытие как граф рождения}
\label{ch:v10-nonequilibrium-bath-universal-cover-growth-graph}

Трёхвершинный ориентированный хопфовский цикл $C_3$ имеет универсальное
покрытие, изоморфное целочисленной прямой $\mathbb Z$. Его проекция
\begin{equation}
 p(n)=n\pmod 3
\end{equation}
переводит каждый ориентированный шаг $n\to n+1$ в следующее ребро цикла.
Шар радиуса три вокруг выбранного подъёма содержит ровно семь вершин
$-3,-2,\ldots,3$ и совпадает с контрольным путём причинного ядра.

Пусть $B_{\rm cov}$ и $B_{C_3}$ --- граничные матрицы пути и цикла, а
$p_0,p_1$ --- проекции вершин и рёбер. Прямой расчёт даёт цепное равенство
\begin{equation}
 B_{C_3}p_1=p_0B_{\rm cov}.
\end{equation}
Обе проекции сюръективны. На пяти внутренних вершинах конечного окна также
выполнено переплетение операторов соседства
\begin{equation}
 p_0A_{\rm cov}=A_{C_3}p_0,
\end{equation}
а сдвиг на три вершины не меняет $p_0$. Концы семиузлового пути являются
только границей выбранного конечного окна, а не концами полного покрытия.

Универсальное покрытие несёт целочисленную функцию высоты $h(n)=n$. Для
каждого ориентированного ребра
\begin{equation}
 B_{\rm cov}^{\top}h=\mathbf 1_6.
\end{equation}
Тем самым один шаг вперёд типизирован как один дискретный акт рождения.
Числа вершин в шарах радиусов $0,1,2,3$ равны $(1,3,5,7)$, оболочки имеют
мощности $(1,2,2,2)$, а будущая полупрямая содержит ровно одну новую вершину
на каждом шаге. Это линейная однопоточная история, а не ветвящееся дерево.

При этом $h$ определена лишь с точностью до $h\mapsto h+c$. Комбинаторное
покрытие не задаёт длительность шага и не добавляет нового размерного
уравнения. После $v_g=c$ прежняя метрическая карта сохраняет ранг/ядро
$8/1$ с той же общей длино-энергетической модой.

\begin{theorem}[Канонический подъём локального пути]
Семиузловой причинный путь является точным конечным окном универсального
покрытия хопфовского цикла. Проекции образуют цепной морфизм, локально
переплетают соседство и переводят поступательный шаг покрытия в циклический
ход основания. Функция высоты даёт один дискретный акт рождения на ребро.
\end{theorem}

\begin{nogocriterion}[Комбинаторная высота не является физическим временем]
Каноничность универсального покрытия не доказывает, что оно является
физическим графом пространства-времени. Оно не выбирает начало высоты,
длительность одного рождения или длину ребра. Эти три утверждения требуют
отдельного физического морфизма, а не только теории покрытий.
\end{nogocriterion}

\begin{protocol}\begin{enumerate}
\item условная архитектура покрытия: \textbf{$12/12$};
\item типизированный цепной морфизм: \textbf{$1/1$};
\item дискретная высота рождения: \textbf{$1/1$};
\item происхождение физического графа роста: \textbf{$0/1$};
\item происхождение абсолютной метрики ребра: \textbf{$0/1$};
\item происхождение физического времени рождения: \textbf{$0/1$};
\item следующий гейт: \textbf{морфизм высоты покрытия в физическое время}.
\end{enumerate}\end{protocol}

Машинный сертификат сохранён в
\path{s2t_v10_cell_birth_four_volume_nonequilibrium_bath_universal_cover_growth_graph_typed_embedding_gate_results.json}.